\subsection{\texorpdfstring{The relation \(\sum_{i} e_{i} f_{i} = n\)}{The relation Σ ei fi = n}}

PROPOSITION 5. Let K be a field, v a valuation on K, A its ring, m its ideal, L a finite extension of K of degree n, B the integral closure of A in L and \((v_{i}^{\prime})_{1 \leq i \leq s}\) a complete system of extensions of v to L. Then

\[
[ \mathrm{B} / \mathrm{mB}: \mathrm{A} / \mathrm{m} ] = \sum_ {i = 1} ^ {s} \varepsilon (v _ {i} ^ {\prime} / v) f (v _ {i} ^ {\prime} / v).
\]

Let \(A_{i}\) be the ring of \(v_{i}^{\prime}\) ; then \(A_{i} = B_{m_{i}}\) , where \(m_{i}\) runs through the family of maximal ideals of B (no. 3, Remark). Let \(q_{i}\) be the saturation of mB with respect to \(m_{i}\) (Chapter II, § 2, no. 4). By Chapter V, Corollary 3 to Proposition 1, no. 1, § 2, the canonical homomorphism \(B/mB \rightarrow \prod_{i=1}^{s} B/q_i\) is an isomorphism and \(m_i\) is the only maximal ideal of B containing \(q_i\) . Hence \(B/q_i\) is canonically isomorphic to \((\mathrm{B}/q_i)_{m_i}\) (Chapter II, §3, no. 3, Proposition 8), that is to

\[
\mathrm{B} _ {\mathrm{m} _ {i}} / \mathrm{mB} _ {\mathrm{m} _ {i}} = \mathrm{A} _ {i} / \mathrm{mA} _ {i}.
\]

Therefore there is a canonical isomorphism \(B/mB \rightarrow \prod_{i=1}^{s} A_{i}/mA_{i}\) , whence the result by virtue of Proposition 4 of no. 4.

COROLLARY. With the same hypotheses and notation,

\[
[ \mathrm{B} / \mathrm{mB}: \mathrm{A} / \mathrm{m} ] = \sum_ {i = 1} ^ {s} \varepsilon (v _ {i} ^ {\prime} / v) f (v _ {i} ^ {\prime} / v) \leqslant \sum_ {i = 1} ^ {s} e (v _ {i} ^ {\prime} / v) f (v _ {i} ^ {\prime} / v) \leqslant n.
\]

We know that \(\varepsilon(v_i'/v) \leqslant e(v_i'/v)\) (no. 4, Corollary to Proposition 3) and \(\sum_{i=1}^{s} e(v_i'/v)f(v_i'/v) \leqslant n\) (no. 3, Theorem 1).

THEOREM 2. With the hypotheses and notation of Proposition 5, the following conditions are equivalent:

\begin{enumerate}
\def\labelenumi{(\alph{enumi})}
\item
  B is a finitely generated A-module;
\item
  B is a free A-module;
\item
  \([\mathbf{B} / \mathfrak{m}\mathbf{B}\colon \mathbf{A} / \mathfrak{m}] = n;\)
\end{enumerate}

\[
\text {(d)} \sum_ {i = 1} ^ {n} e (v _ {i} ^ {\prime} / v) f (v _ {i} ^ {\prime} / v) = n \text { and } \varepsilon (v _ {i} ^ {\prime} / v) = e (v _ {i} ^ {\prime} / v) \text { for   all } i.
\]

The equivalence of (a) and (b) follows from Lemma 1, § 3, no. 6. Clearly (b) implies (c) (Algebra, Chapter II, § 1, no. 5, formula (19)). The equivalence of (c) and (d) follows from the Corollary to Proposition 5. It remains to show that (c) implies (b).

In general, if M is an A-module, we shall denote by V(M) the vector space M/mM over A/m. Hypothesis (c) means that \(\dim(\mathrm{V}(\mathrm{B})) = n\) . Let \(x_{1}, \ldots, x_{n}\) be elements of B whose canonical images in V(B) form a basis of V(B) and let \(L \subset B\) be the sub-A-module which they generate. As L is torsion-free and finitely generated, it is free (§ 3, no. 6, Lemma 1). We shall see that B = L. Let \(y \in B\) ; we write \(M = L + Ay\) ; this is also a free A-module. The canonical injections \(L \to M \to B\) give canonical homomorphisms \(\mathrm{V}(\mathrm{L}) \to \mathrm{V}(\mathrm{M}) \to \mathrm{V}(\mathrm{B})\) . As the ranks of L and M are \(\leqslant n\) , so are the dimensions of V(L) and V(M). Now, by hypothesis, \(\mathrm{V}(\mathrm{L}) \to \mathrm{V}(\mathrm{B})\) is surjective and V(B) is n-dimensional hence V(L) and V(M) are n-dimensional and \(\mathrm{V}(\mathrm{L}) \to \mathrm{V}(\mathrm{M})\) is surjective. As M is finitely generated, \(L \to M\) is surjective (Chapter II, § 3, no. 2, Corollary 1 to Proposition 4), whence \(L = M, y \in L\) and B = L. Hence B is free.

Remark (1) If \(v\) is discrete, \(\varepsilon(v_i'/v) = e(v_i'/v)\) (no. 4) and condition (d) reduces to \(\sum_{i=1}^{s} e(v_i'/v)f(v_i'/v) = n\) .

COROLLARY 1. With the same hypotheses and notation, suppose further that v is discrete and L separable. Then

\[
\sum_ {i = 1} ^ {s} e (v _ {i} ^ {\prime} / v) f (v _ {i} ^ {\prime} / v) = n.
\]

The integral closure B of A is then a free A-module of rank n, since A is a principal ideal domain (Chapter V, § 1, no. 6, Corollary 2 to Proposition 18).

COROLLARY 2. Let K be a field, v a discrete valuation on K with respect to which K is complete and L a finite extension of K of degree n.~Then v admits a unique (up to equivalence) extension v' to L, the ring A' of v' is a finitely generated free module over the ring A of v and \(e(v'/v)f(v'/v)=n\) .

All the extensions of \(v\) to \(L\) are dependent (no. 2, Corollary to Proposition 2); since they are discrete (no. 1, Corollary 3 to Proposition 1), they are therefore equivalent (§ 4, no. 5, Proposition 6 (c)). This shows the uniqueness of \(v'\) . The integral closure of \(A\) in \(L\) is therefore \(A'\) (§ 1, no. 3, Corollary 3 to Theorem 3). As \(v\) is discrete, the topology induced on \(A\) by that on \(K\) is the m-adic topology (where \(m = m(A)\) ); the ring \(A\) is complete, for it is closed in \(K\) . We conclude that, since \(A'/mA'\) is a finite-dimensional vector \((A/m)\) -space (no. 4, Proposition 4), \(A'\) is a finitely generated A-module (Chapter III, § 2, no. 9, Corollary 3 to Proposition 12). It is therefore free and \(e(v'/v)f(v'/v) = n\) by Theorem 2.

COROLLARY 3. Suppose that \(v\) is of height 1 and that the equivalent conditions of Theorem 2 hold; if \(\hat{\mathbf{L}}_i\) is the completion of \(\mathbf{L}\) with respect to \(v_i'\) , the degree \(n_i = [\hat{\mathbf{L}}_i: \hat{\mathbf{K}}]\) is equal to \(e(v_i'|v)f(v_i'|v)\) for all \(i\) and the canonical homomorphism

\[
\phi \colon \hat {\mathbf {K}} \otimes_ {\mathbf {K}} \mathbf {L} \rightarrow \prod_ {i = 1} ^ {s} \hat {\mathbf {L}} _ {i}
\]

(no. 2, Proposition 2). is bijective. For all \(x \in \mathbf{L}\) , the characteristic polynomial \(\mathrm{Pc}_{\mathrm{L} / \mathbb{K}}(x; \mathbf{X})\) is equal to the product of the characteristic polynomials \(\mathrm{Pc}_{\hat{\mathrm{L}}_i / \hat{\mathbf{K}}}(x; \mathbf{X})\) ( \(1 \leqslant i \leqslant s\) ); in particular,

\[
\left\{ \begin{array} { l } \mathrm{Tr} _ { \mathbf { L } / \mathbb { K } } ( x ) = \sum _ { i = 1 } ^ { s } \mathrm{Tr} _ { \hat { \mathbf { L } } _ { i } / \hat { \mathbf { K } } } ( x ) \\ \mathrm{N} _ { \mathbf { L } / \mathbb { K } } ( x ) = \prod _ { i = 1 } ^ { s } \mathrm{N} _ { \hat { \mathbf { L } } _ { i } / \hat { \mathbf { K } } } ( x ) \\ v ( \mathrm{N} _ { \mathbf { L } / \mathbb { K } } ( x ) ) = \sum _ { i = 1 } ^ { s } n _ { i } v _ { i } ^ { \prime } ( x ) . \\ \end{array} \right.\tag{9}
\]

(The last relation in (9) is meaningful, for we may obviously assume that the \(v_{i}^{\prime}\) , which are of height 1 by Corollary 2 to Proposition 1 of no. 1, take, as does \(v\) , their values in a subgroup of \(\mathbf{R}\) .)

As no two of the \(v_{i}^{\prime}\) are equivalent and they are of height 1, they are independent and Proposition 2 of no. 2 shows therefore that \(e(v_{i}^{\prime}/v)f(v_{i}^{\prime}/v) \leqslant n_{i}\) for all \(i\) and \(\sum_{i=1}^{s} n_i \leqslant n\) . The first assertion therefore follows from these inequalities and the relation \(\sum_{i=1}^{s} e(v_i'/v) f(v_i'/v) = n\) . Under the isomorphism \(\phi\) the endomorphism \(z \mapsto z(1 \otimes x)\) of \(\hat{\mathbf{K}} \otimes_{\mathbb{K}} \mathbf{L}\) (for \(x \in \mathbf{L}\) ) is transformed into the endomorphism of \(\prod_{i=1}^{s}\hat{L}_i\) leaving invariant each of the factors and reducing on each factor to multiplication by \(x\) (L being canonically imbedded in its completion \(\hat{L}_i\) ); whence the assertion relating to the characteristic polynomial of \(x\) and the first two formulae of (9). Finally, let E be a finite quasi-Galois extension of \(\hat{K}\) , containing \(\hat{L}_i\) ; as \(\hat{K}\) is complete and \(\hat{v}\) of height 1, there exists only one valuation (up to equivalence) \(w\) on E extending \(\hat{v}\) (no. 2, Corollary 1 to Proposition 2); then, for every \(\hat{K}\) -automorphism \(\sigma\) of E, \(w(\sigma(x)) = v_i'(x)\) . Therefore

\[
\hat {v} \left(\mathrm{N} _ {\hat {\mathrm{L}} _ {i} / \hat {\mathrm{K}}} (x)\right) = n _ {i} v _ {i} ^ {\prime} (x)
\]

(Algebra, Chapter VIII, § 12, no. 2, formula (15)), which proves the formula of (9).

COROLLARY 4. Under the hypotheses of Corollary 3, if L is a separable extension of K, each of the \(\hat{\mathbf{L}}_i\) is a separable extension of \(\hat{\mathbf{K}}\) . If further L is a Galois extension of K with Galois group \(\mathcal{G}\) and \(\mathcal{G}_i\) denotes the decomposition group of the ideal of \(v_i'\) in B (Chapter V, § 2, no. 2, Definition 2), then \(\hat{\mathbf{L}}_i\) is a Galois extension of \(\hat{\mathbf{K}}\) whose Galois group is isomorphic to \(\mathcal{G}_i\) .

Clearly \(\hat{\mathbf{L}}_i = \hat{\mathbf{K}}(\mathbf{L})\) ; hence, if \(\mathbf{L}\) is separable over \(\mathbf{K}\) , \(\hat{\mathbf{L}}_i\) is separable over \(\hat{\mathbf{K}}\) (Algebra, Chapter V, § 7, no. 6, Proposition 10). Suppose now that \(\mathbf{L}\) is Galois. Every automorphism \(\sigma \in \mathcal{G}_i\) is continuous on \(\mathbf{L}\) with the topology defined by \(v_i'\) , the fact that no two of the ideals of the \(v_i'\) are comparable with respect to inclusion (§ 7, no. 2, Corollary 1 to Theorem 1) necessarily implying that \(v_i' = v_i' \circ \sigma\) by definition of \(\mathcal{G}_i\) ; hence \(\sigma\) may be extended by continuity to a \(\hat{\mathbf{K}}\) -automorphism \(\hat{\sigma}\) of \(\hat{\mathbf{L}}_i\) . This proves that the number of \(\hat{\mathbf{K}}\) -automorphisms of \(\hat{\mathbf{L}}_i\) is at least equal to \(\operatorname{Card}(\mathcal{G}_i)\) . But as the valuations \(v_i'\) are pairwise conjugate under \(\mathcal{G}\) (Chapter V, § 2, no. 3, Proposition 6), \(s = (\mathcal{G}: \mathcal{G}_i)\) , whence

\[
\operatorname{Card} \left(\mathcal {G} _ {i}\right) = n / s \leqslant n.
\]

and on the other hand \(n = sn_{i}\) by Corollary 3; this proves that \(\hat{\mathbf{L}}_{i}\) is a Galois extension of \(\hat{\mathbf{K}}\) and that the extensions by continuity of the automorphisms \(\sigma \in \mathcal{G}_{i}\) are the only \(\hat{\mathbf{K}}\) -automorphisms of \(\hat{\mathbf{L}}_{i}\) .

Remark (2). Part of the above results extends to the case of valuations on a not necessarily commutative field K (cf.~§ 3, no. 1). Let L be an extension field of K and let \(v'\) be a valuation on L, \(v\) its restriction to K and A' and A the respective rings of the valuations \(v'\) and v; then there is defined a ramification index \(e(v'/v)\) as in no. 1; on the other hand, \(\kappa(\mathbf{A})\) is identified with a subfield of \(\kappa(\mathbf{A}')\) and the (left) residue rank of \(v'\) with respect to v is defined to be the number \(f(v'/v)\) equal to the dimension of the left vector \(\kappa(\mathbf{A})\) -space \(\kappa(\mathbf{A}')\) , if this dimension is finite, and \(+\infty\) in the opposite case. Then, if L is a left vector K-space of finite dimension n, Lemma 2 of no. 1 and its proof go over unchanged. Moreover, if K is complete with respect to v, the assertions of Corollary 2 to Theorem 2 of no. 5 (other than the existence of \(v'\) ) are also valid (n denoting the dimension of L as a left vector K-space) with the following proof:

In the first place the topology defined by \(v'\) on \(L\) is Hausdorff and compatible with its left vector \(K\) -space structure and hence two extensions of \(v\) to \(L\) give the same topology on \(L\) (§ 5, no. 2, Proposition 4), which proves that these extensions are the same up to equivalence (§ 6, no. 2). We show next that, if \(m = m(A)\) , \(A'/mA'\) is a left vector \((A/m)\) -space of dimension \(e(v'/v)f(v'/v)\) . Writing \(e = e(v'/v)\) , we may assume that \(v(K^*) = Z\) and \(v'(L^*) = e^{-1}Z\) ; let \(u'\) be an element of \(L\) such that \(v'(u') = e^{-1}\) and \(u\) an element of \(K\) such that \(v(u) = 1\) ; hence \(u = zu'^e\) , where \(z \in L\) is such that \(v'(z) = 0\) . As \(m\) is generated by \(u\) (as a left or right ideal of \(A\) ), \(mA' = u'^eA' = A'u'^e\) and it suffices to prove that, for \(0 \leqslant k \leqslant e - 1\) , \(A'u'^k/A'u'^{k+1}\) is a left vector \((A/m)\) -space of dimension \(f(v'/v)\) . But \(t \mapsto tu'^k\) is an isomorphism of the left A-module \(A'\) onto the left A-module \(A'u'^k\) mapping \(A'u'\) to \(A'u'^{k+1}\) and which therefore gives by taking quotients an \((A/m)\) -isomorphism of \(A'/A'u'\) onto \(A'u'^k/A'u'^{k+1}\) , whence our assertion by definition of \(f(v'/v)\) , \(u'\) generating the maximal ideal of \(A'\) . The proof is completed as when \(K\) and \(L\) are commutative (the fact that a finitely generated torsion-free A-module is free being proved as in § 3, no. 6, Lemma 1).

